\section*{Problem 7 --- Solution}

\partlabel{(a)}
For a single-phase, fixed-composition substance, let $h$, $s$, and $v$ be
mass-specific properties. Since $h=h(T,p)$,
% L2 notes, PDF p.6 (printed p.21), Eq. (40): cp definition.
\begin{align*}
dh&=\left(\frac{\partial h}{\partial T}\right)_p dT
   +\left(\frac{\partial h}{\partial p}\right)_T dp
 =c_p\,dT+\left(\frac{\partial h}{\partial p}\right)_T dp.
\end{align*}
% L3 notes, PDF p.4 (printed p.28), Eqs. (91), (97)--(99).
Using $dh=T\,ds+v\,dp$ and the mass-specific Gibbs energy $g=h-Ts$,
\begin{align*}
dg&=dh-T\,ds-s\,dT
   =(T\,ds+v\,dp)-T\,ds-s\,dT=-s\,dT+v\,dp.
\end{align*}
\[
\left(\frac{\partial g}{\partial T}\right)_p=-s,\qquad
\left(\frac{\partial g}{\partial p}\right)_T=v.
\]
The mixed derivatives of the state function $g$ are equal, so
\[
\frac{\partial}{\partial p}\left(\frac{\partial g}{\partial T}\right)_p
=\frac{\partial}{\partial T}\left(\frac{\partial g}{\partial p}\right)_T
\quad\Longrightarrow\quad
\left(\frac{\partial s}{\partial p}\right)_T
=-\left(\frac{\partial v}{\partial T}\right)_p.
\]
Therefore, at fixed $T$,
\[
\left(\frac{\partial h}{\partial p}\right)_T
=T\left(\frac{\partial s}{\partial p}\right)_T+v
=v-T\left(\frac{\partial v}{\partial T}\right)_p,
\]
% L3 notes, PDF p.5 (printed p.29), Eq. (106).
and substitution gives
\[
\boxed{dh=c_p\,dT+
\left[v-T\left(\frac{\partial v}{\partial T}\right)_p\right]dp.}
\]

\partlabel{(b)}
% L3 notes, PDF p.2 (printed p.26), Eq. (79), and PDF p.6,
% printed p.30, Eqs. (112)--(113).
With the mass-specific gas constant $R$, $Z=pv/(RT)$ gives
\[
v=\frac{ZRT}{p},\qquad
\left(\frac{\partial v}{\partial T}\right)_p
=\frac{R}{p}\left[Z+T\left(\frac{\partial Z}{\partial T}\right)_p\right].
\]
Thus,
\begin{align*}
v-T\left(\frac{\partial v}{\partial T}\right)_p
&=\frac{ZRT}{p}-\frac{RT}{p}
  \left[Z+T\left(\frac{\partial Z}{\partial T}\right)_p\right]\\
&=\frac{ZRT}{p}-\frac{ZRT}{p}
  -\frac{RT^2}{p}\left(\frac{\partial Z}{\partial T}\right)_p
=-\frac{RT^2}{p}\left(\frac{\partial Z}{\partial T}\right)_p.
\end{align*}
Hence
\[
\boxed{dh=c_p\,dT-\frac{RT^2}{p}
\left(\frac{\partial Z}{\partial T}\right)_p dp.}
\]
For an ideal gas, $Z=1$, so the pressure term vanishes.

\newpage
\partlabel{(c)}
The supplied $B$ is molar; use $\bar v$, $\bar h$, and
$R_u=8.314462618\ \mathrm{J\,mol^{-1}K^{-1}}$. Also,
\[
\mathcal M=0.0280134\ \mathrm{kg\,mol^{-1}},\qquad
\Delta p=(200-1)\times10^5=1.99\times10^7\ \mathrm{Pa}.
\]
% HW1 PDF p.3, P7(c): model and data. L3 notes, PDF p.11
% (printed p.35), Eqs. (140)--(141): corresponding molar model.
Within the supplied pressure-truncated model,
\begin{align*}
Z&=1+\frac{B(T)p}{R_uT},
&\bar v&=\frac{ZR_uT}{p}=\frac{R_uT}{p}+B(T),\\
\left(\frac{\partial\bar v}{\partial T}\right)_p
&=\frac{R_u}{p}+B'(T),
&B'(T)&\equiv\frac{dB}{dT}.
\end{align*}
Since $dT=0$, the molar form of part (a) gives
\begin{align*}
d\bar h
&=\left[\frac{R_uT}{p}+B
-T\left(\frac{R_u}{p}+B'\right)\right]dp
=(B-TB')\,dp,\\
\Delta\bar h&=\int_{p_1}^{p_2}(B-TB')\,dp=(B-TB')\Delta p.
\end{align*}
\[
\boxed{\Delta h=\frac{\Delta\bar h}{\mathcal M}
=\frac{(B-TB')\Delta p}{\mathcal M}.}
\]

\textbf{(i) Rigid spheres.}
With $N_A=6.02214076\times10^{23}\ \mathrm{mol^{-1}}$,
% HW1 PDF p.3, P7(c)(i), supplies B_RS; L3 notes PDF p.11,
% Eqs. (137)--(138), give the constant-excluded-volume limit.
\begin{align*}
B_{\rm RS}
&=\frac{2\pi}{3}N_A\sigma^3
=\frac{2\pi}{3}(6.02214076\times10^{23})(3.681\times10^{-10})^3\\
&=6.2908158\times10^{-5}\ \mathrm{m^3\,mol^{-1}},
\qquad B'_{\rm RS}=0,\\
\Delta\bar h_{\rm RS}
&=(6.2908158\times10^{-5})(1.99\times10^7)
=1251.872\ \mathrm{J\,mol^{-1}},\\
\Delta h_{\rm RS}
&=\frac{1251.872}{0.0280134}\ \mathrm{J\,kg^{-1}}
=\boxed{+44.69\ \mathrm{kJ\,kg^{-1}}}.
\end{align*}

\textbf{(ii) Attraction-inclusive model.}
% HW1 PDF p.3, P7(c)(ii): B and its temperature derivative.
\begin{align*}
B-TB'&=[-75.7-150(1.01)]\times10^{-6}
=-227.2\times10^{-6}\ \mathrm{m^3\,mol^{-1}},\\
\Delta\bar h&=(-227.2\times10^{-6})(1.99\times10^7)
=-4521.280\ \mathrm{J\,mol^{-1}},\\
\Delta h&=\frac{-4521.280}{0.0280134}\ \mathrm{J\,kg^{-1}}
=\boxed{-161.4\ \mathrm{kJ\,kg^{-1}}}.
\end{align*}
% L3 notes, PDF p.11 (printed p.35), Example D;
% enthalpy definition also PDF p.15 (printed p.39), Eq. (173).
At fixed $T$, the same model gives
\[
\Delta\bar e=\Delta\bar h-\Delta(p\bar v)
=(B-TB')\Delta p-B\Delta p=-TB'\Delta p.
\]
Rigid spheres have $B'=0$: internal energy is unchanged, and the positive
enthalpy change comes from $p\bar v$. Attractions introduce a negative
internal-energy contribution; here $B-TB'<0$, reversing the enthalpy trend.

% L3 notes, PDF p.9 (printed p.33), Eqs. (125)--(126):
% the virial expansion is a low-density expansion.
At $200\ \mathrm{bar}$, however,
\[
Z_{\rm RS}=2.0088,\qquad Z_{\rm attraction}=-0.21395.
\]
The rigid-sphere correction is not small, and the negative $Z$ in model (ii)
is unphysical. These are formal results of the assigned truncation;
model (ii) fails before the final pressure is reached.